\documentclass[10pt, aspectratio=169]{beamer}

\usepackage{fontspec}
\usepackage[spanish,es-noshorthands]{babel}
\usepackage{../plantilla/latex/unab_presentacion}

% Los datos compartidos se generan desde datos/datos.md al ejecutar compilar.py.
\input{../.build/datos.tex}
\input{../.build/citas.tex}
\title{\TituloTrabajo}
\author{\Grupo}
\hypersetup{pdftitle={\TituloTrabajo},pdfauthor={\Grupo}}
\setdatosfooter{\TituloCorto}{\Asignatura}

\begin{document}

% =========================================================
% Lámina 1: Portada
% =========================================================
\portadainstitucional

% =========================================================
% Lámina 2: Reflexión
% =========================================================
\expositor{\Grupo}
\begin{frame}{\faLightbulb \ Reflexión Inicial}
  \begin{center}
    \vfill
    \IfFileExists{../assets/contenido/imagenes/reflexion.png}{%
      \includegraphics[width=0.75\linewidth, height=0.75\textheight, keepaspectratio]{../assets/contenido/imagenes/reflexion.png}
    }{%
      \makebox[\textwidth][c]{%
        \begin{beamercolorbox}[rounded=true,wd=0.78\textwidth,sep=0.45cm,center]{block body}
          \Large\textcolor{unabblue}{\faLightbulb}\\[0.25cm]
          \large ¿Qué problema queremos resolver y por qué importa?
        \end{beamercolorbox}%
      }
    }
    \vfill
  \end{center}
\end{frame}

% =========================================================
% Lámina 3: Conceptos
% =========================================================
\expositor{\Grupo}
\begin{frame}{\faInfoCircle \ Conceptos Fundamentales}
  \begin{block}{Concepto principal}
    Escribe aquí la definición central del tema. Mantén una idea clara, breve y orientada a la exposición oral. Ejemplo de cita: \cita{fuente}.
  \end{block}

  \begin{columns}[onlytextwidth]
    \column{0.48\textwidth}
    \begin{itemize}
      \item \textbf{\faCheckCircle \ Punto 1:} descripción breve.
      \item \textbf{\faCheckCircle \ Punto 2:} descripción breve.
    \end{itemize}
    \column{0.48\textwidth}
    \begin{itemize}
      \item \textbf{\faCheckCircle \ Punto 3:} descripción breve.
      \item \textbf{\faCheckCircle \ Punto 4:} descripción breve.
    \end{itemize}
  \end{columns}
\end{frame}

% =========================================================
% Lámina 4: Desarrollo
% (Nota: si agregas código fuente con lstlisting en un frame, usa \begin{frame}[fragile]{Título})
% =========================================================
\expositor{\Grupo}
\begin{frame}{\faMicrochip \ Desarrollo del Tema}
  \begin{block}{Idea operativa}
    Explica aquí cómo funciona el tema, qué componentes intervienen y por qué es importante.
  \end{block}

  \begin{columns}[T]
    \column{0.48\textwidth}
    \begin{exampleblock}{Subtema A}
      \begin{itemize}
        \item Característica principal.
        \item Ventaja o aplicación.
        \item Ejemplo breve.
      \end{itemize}
    \end{exampleblock}
    \column{0.48\textwidth}
    \begin{exampleblock}{Subtema B}
      \begin{itemize}
        \item Característica principal.
        \item Ventaja o aplicación.
        \item Ejemplo breve.
      \end{itemize}
    \end{exampleblock}
  \end{columns}
\end{frame}

% =========================================================
% Lámina 5: Casos de uso
% =========================================================
\expositor{\Grupo}
\begin{frame}{\faBriefcase \ Aplicaciones y Casos de Uso}
  \begin{block}{Dónde se aplica}
    \begin{itemize}
      \item \textbf{\faBriefcase \ Caso 1:} descripción del caso de uso.
      \item \textbf{\faCogs \ Caso 2:} descripción del caso de uso.
      \item \textbf{\faUsers \ Caso 3:} descripción del caso de uso.
      \item \textbf{\faChartLine \ Caso 4:} descripción del caso de uso.
    \end{itemize}
  \end{block}

  \begin{center}
    \textcolor{unabblue}{\Large \faLightbulb \quad \faProjectDiagram \quad \faClipboardCheck \quad \faChartBar}
  \end{center}
\end{frame}

% =========================================================
% Lámina 6: Comparativa
% =========================================================
\expositor{\Grupo}
\begin{frame}{\faBalanceScale \ Análisis Comparativo}
  \begin{columns}[T, onlytextwidth]
    \column{0.48\textwidth}
    \begin{block}{\faArrowLeft \ Enfoque A}
      \begin{itemize}
        \item \textbf{Criterio 1:} descripción.
        \item \textbf{Criterio 2:} descripción.
        \item \textbf{Criterio 3:} descripción.
        \item \textbf{Criterio 4:} descripción.
      \end{itemize}
    \end{block}

    \column{0.48\textwidth}
    \begin{block}{\faArrowRight \ Enfoque B}
      \begin{itemize}
        \item \textbf{Criterio 1:} descripción.
        \item \textbf{Criterio 2:} descripción.
        \item \textbf{Criterio 3:} descripción.
        \item \textbf{Criterio 4:} descripción.
      \end{itemize}
    \end{block}
  \end{columns}

  \vspace{0.35cm}
  \begin{center}
    \begin{beamercolorbox}[rounded=true, wd=0.8\textwidth, sep=0.12cm, center]{block body}
      \footnotesize \textcolor{unabblue}{\faLightbulb \ \textbf{Punto clave:}} escribe aquí la conclusión de la comparativa.
    \end{beamercolorbox}
  \end{center}
\end{frame}

% =========================================================
% Lámina 7: Conclusiones
% =========================================================
\expositor{\Grupo}
\begin{frame}{\faCheckCircle \ Conclusiones Estratégicas}
  \begin{block}{Idea final}
    Resume aquí el mensaje más importante de la exposición.
  \end{block}

  \begin{columns}[T]
    \column{0.48\textwidth}
    \begin{exampleblock}{\faBullhorn \ Conclusión 1}
      \footnotesize
      \begin{itemize}
        \item Argumento de cierre.
        \item Beneficio principal.
        \item Aplicación práctica.
      \end{itemize}
    \end{exampleblock}

    \column{0.48\textwidth}
    \begin{exampleblock}{\faTools \ Conclusión 2}
      \footnotesize
      \begin{itemize}
        \item Argumento de cierre.
        \item Beneficio principal.
        \item Aplicación práctica.
      \end{itemize}
    \end{exampleblock}
  \end{columns}
\end{frame}

% =========================================================
% Lámina 8: Referencias
% =========================================================
\expositor{\Grupo}
\begin{frame}{\faBook \ Referencias Bibliográficas}
  \begin{block}{Fuentes consultadas}
    \justifying
    \footnotesize
    Agrega aquí las fuentes consultadas para respaldar la exposición.
  \end{block}

  \input{../.build/referencias-presentacion.tex}
\end{frame}

% =========================================================
% Lámina 9: Cierre
% =========================================================
\cierreinstitucional[Muchas gracias por su atención]

\end{document}
